\chapter{Sequence Models on Order Books}\label{ch:ml:sequence-models-on-order-books}

A widely cited paper reports that a convolutional-recurrent network reading the raw states of a ten-level order book
outperforms all existing algorithms on a benchmark order-book dataset, and delivers stable out-of-sample accuracy on a
year of London Stock Exchange quotes. On a laptop,
on a few hours of a simulated book, with labels built honestly and every test session kept out of training, the best
sequence model reaches an information coefficient of 0.43, and a \omterm{def:ml:linear-and-regularised-baselines:logistic}{logistic regression} on ten hand-made features (the
order-flow imbalance, the queue imbalance, recent flow and moves) reaches 0.54. The paper is not wrong about its data; the result is that
raw books are a hard input, sequence models need a lot of them, and the claim ``deep learning finds the features by
itself'' has a data cost that must be paid in sessions. This chapter builds the four standard sequence architectures,
trains them on order-book windows, and measures what they buy.

\section{Order-book streams as inputs}

The chapter's data are ten-minute \texttt{firm.tape} sessions (Book~7, chapter~2), each on its own seed: eight for
training, three for \omterm{def:ml:trees-and-boosting:early}{early stopping}, eight for testing, and eight more to double the \omterm{def:ml:why-financial-machine-learning-is-different:sets}{training set} later. Every half
second the book is rebuilt from the messages and summarised level by level: for each of the ten best levels on each side,
the distance of the price from the mid, in ticks, and the logarithm of the size in lots, forty numbers per snapshot. An
input is a window of the last 20 snapshots (ten seconds); its label is the direction of the mean mid over the next ten
seconds against the current mid, down, flat or up, with a dead band of a quarter tick (\cref{fig:ml:lobseq:window}).
Windows never cross sessions, and the scaling of the forty columns is fitted on training sessions only.

\begin{omfigure}
\begin{tikzpicture}[x=0.42cm,y=0.42cm,font=\footnotesize]
\foreach \t in {0,...,19} { \foreach \l in {0,...,9} {
  \fill[omDef!15] (\t,\l) rectangle ++(0.85,0.85); } }
\draw[omDef,thick] (-0.2,-0.2) rectangle (19.95,10.05);
\node[anchor=east,align=right] at (-0.5,5) {10 levels $\times$ 4\\(distance, log size,\\bid and ask)};
\node at (10,-1.1) {20 snapshots, one every half second (ten seconds)};
\draw[->,thick] (20.4,5) -- (23.3,5) node[midway,above] {model};
\node[draw=omThm,fill=omThm!10,rounded corners=2pt,align=center,anchor=west] at (23.5,5) {$\P(\text{down})$, $\P(\text{flat})$,\\ $\P(\text{up})$};
\draw[omProp,thick,->] (19.9,11.2) -- (26.5,11.2);
\node[omProp,anchor=south] at (23.2,11.3) {next 10 s: mean mid $-$ mid};
\draw[omProp,dashed] (19.9,10.2) -- (19.9,11.6);
\node[anchor=south] at (10,10.2) {decision time $t$ at the window's end};
\end{tikzpicture}

\omcaption{An input window and its label. The label compares the average mid over the next ten seconds with the mid at
the decision time and is flat within a quarter tick; the smoothing follows the DeepLOB literature and makes the label
easier to predict than the next move itself.}\label{fig:ml:lobseq:window}
\end{omfigure}

The classes are unbalanced: 73\% of the test windows are flat, 15\% down and 13\% up. Always predicting ``flat'' is
right 73\% of the time, so accuracy is a poor score; the chapter reports it for comparison with the literature and ranks
the models by the information coefficient of $\P(\text{up}) - \P(\text{down})$ with the forward change, session by
session.

\section{Convolutional and recurrent models}

\begin{definition}[Convolutional neural network]\label{def:ml:sequence-models-on-order-books:cnn}
A \emph{convolutional neural network}\index{convolutional neural network} applies the same small filter at every
position of its input: a one-dimensional convolutional layer maps a sequence $z_1,\dots,z_T$ of $d$-vectors to
$u_t = \mathrm{act}(\sum_{j=0}^{k-1}B_jz_{t+j-\lfloor k/2\rfloor} + c)$, with $k$ filter matrices $B_j$ shared across $t$.
\end{definition}

\begin{definition}[Recurrent neural network, long short-term memory]\label{def:ml:sequence-models-on-order-books:rnn}
A \emph{recurrent neural network}\index{recurrent neural network} reads a sequence one step at a time and carries a hidden
state, $h_t = \mathrm{act}(A x_t + C h_{t-1} + c)$. A \emph{long short-term memory}\index{long short-term memory} network
(LSTM) adds a cell state updated through gates (input, forget and output, each a sigmoid of $x_t$ and $h_{t-1}$), which
lets gradients flow over long sequences (Hochreiter and Schmidhuber, 1997).
\end{definition}

The CNN of the chapter has two convolutional layers of 32 filters of width five, averages over the window and classifies;
the LSTM has 32 units and classifies from its last state. Both read the forty columns of every snapshot of the window.

\section{Temporal convolutions and attention}

\begin{definition}[Causal convolution, temporal convolutional network]\label{def:ml:sequence-models-on-order-books:tcn}
A \emph{causal convolution}\index{causal convolution} uses only the current and earlier positions:
$u_t = \sum_{j=0}^{k-1}B_jz_{t - jr}$, with dilation $r$. A \emph{temporal convolutional network}\index{temporal
convolutional network} (TCN) stacks causal convolutions with dilations $1, 2, 4,\dots$, so that the receptive field grows
exponentially with depth, usually with residual connections (Bai, Kolter and Koltun, 2018).
\end{definition}

\begin{definition}[Attention mechanism, transformer]\label{def:ml:sequence-models-on-order-books:attention}
An \emph{attention mechanism}\index{attention mechanism} lets each position of a sequence form its output as a weighted
average of all positions: with queries $\mathsf Q = ZA_{\mathsf Q}$, keys $\mathsf K = ZA_{\mathsf K}$ and values
$\mathsf V = ZA_{\mathsf V}$ computed from the sequence $Z$ ($T\times h$), the output is
$\mathrm{softmax}(\mathsf Q\mathsf K^\top/\sqrt h)\mathsf V$. A \emph{transformer}\index{transformer} layer combines
several attention heads, a position-wise network, residual connections and \omterm{def:ml:neural-networks-for-noisy-tabular-data:dropout}{layer normalisation}; positions are encoded
by adding learned or fixed vectors to the inputs (Vaswani et al., 2017).
\end{definition}

The TCN has four causal blocks (dilations 1 to 8, filters of width three, receptive field 31 snapshots, more than the
window) and classifies from the last position; the attention model is one \omterm{def:ml:sequence-models-on-order-books:attention}{transformer} layer with two heads of width 32
and learned positions, also read at the last position.

\begin{omfigure}
\begin{tikzpicture}
\begin{axis}[width=10.5cm,height=5.4cm,ybar,bar width=14pt,ylabel={test information coefficient},
  symbolic x coords={logistic,CNN,LSTM,TCN,attention},xtick=data,
  xticklabels={{logistic on\\features},CNN,LSTM,TCN,attention},xticklabel style={align=center},tick label style={font=\small},
  label style={font=\small},ymin=-0.2,ymax=0.7,grid=major,grid style={gray!25},enlarge x limits=0.14,
  table/col sep=comma]
\addplot[fill=omDef!55,draw=omDef,error bars/.cd,y dir=both,y explicit]
  table[x=model,y=ic,y error=sd]{figdata/ml/08-sequence-models-on-order-books/models.csv};
\end{axis}
\end{tikzpicture}

\omcaption{Information coefficient of $\P(\text{up}) - \P(\text{down})$ with the forward change, averaged over eight test
sessions (and two seeds for the sequence models); bars show the standard deviation across sessions and seeds. The
\omterm{def:ml:linear-and-regularised-baselines:logistic}{logistic regression} reads ten hand-made features at the decision time; the sequence models read the raw ten-level book
over ten seconds. Data: \texttt{ml\_lobseq.table}.}\label{fig:ml:lobseq:models}
\end{omfigure}

\section{What the literature replicates}

On the chapter's data (\cref{fig:ml:lobseq:models}), the multinomial \omterm{def:ml:linear-and-regularised-baselines:logistic}{logistic regression} on the hand-built features of
Book~8 (chapter~14; order-flow imbalance over three windows, queue imbalance, flow and past moves) reaches an accuracy of
0.781 and an IC of 0.543. Among the sequence models the \omterm{def:ml:sequence-models-on-order-books:attention}{transformer} does best (accuracy 0.775, IC 0.426), then the TCN
(0.766, 0.366), the LSTM (0.729, 0.292) and the CNN (0.727, 0.060): the CNN averages its filters over the window and so
cannot tell the last snapshot from the first, and it barely beats always predicting ``flat''. The spread across sessions
and seeds (0.15 to 0.20 of IC for the sequence models, 0.09 for the \omterm{def:ml:linear-and-regularised-baselines:logistic}{logistic regression}) is as large as the gaps.

The models were early-stopped after one to four \omterm{def:ml:neural-networks-for-noisy-tabular-data:backprop}{epochs}: eight sessions are 8\,808 windows, and a model with thousands
of parameters that reads 800 numbers per window fits their noise quickly. Doubling the data helps (\cref{fig:ml:lobseq:scaling}):
the TCN's IC goes from 0.13 on four sessions to 0.30 on eight and 0.44 on sixteen, while the \omterm{def:ml:linear-and-regularised-baselines:logistic}{logistic regression} stays
between 0.54 and 0.57 throughout. The sequence models are learning from the raw book what the hand-made features give the
linear model directly, and they are still learning it at sixteen sessions.

\begin{omfigure}
\begin{tikzpicture}
\begin{axis}[width=9.5cm,height=5.2cm,xlabel={training sessions (ten minutes each)},ylabel={test IC},
  tick label style={font=\small},label style={font=\small},xmin=2,xmax=18,ymin=0,ymax=0.7,xtick={4,8,16},
  grid=major,grid style={gray!25},legend columns=2,legend style={font=\footnotesize,at={(0.5,-0.3)},anchor=north,
  draw=none,fill=none,/tikz/every even column/.append style={column sep=8pt}},table/col sep=comma]
\addplot[omDef,very thick,mark=*] table[x=sessions,y=tcn]{figdata/ml/08-sequence-models-on-order-books/scaling.csv};
\addplot[omThm,very thick,mark=square*] table[x=sessions,y=logistic]{figdata/ml/08-sequence-models-on-order-books/scaling.csv};
\legend{TCN on the raw book,logistic on hand-made features}
\end{axis}
\end{tikzpicture}

\omcaption{Test IC as the \omterm{def:ml:why-financial-machine-learning-is-different:sets}{training set} grows from four to sixteen sessions (one seed; the same eight test sessions). Data:
\texttt{ml\_lobseq.scaling}.}\label{fig:ml:lobseq:scaling}
\end{omfigure}

The published evidence points the same way once the evaluation is strict. Zhang, Zohren and Roberts's DeepLOB outperformed
the earlier algorithms on the FI-2010 benchmark (Ntakaris et al.) and kept a stable out-of-sample accuracy on a year of
London Stock Exchange quotes, including instruments it was not trained on. Sirignano and Cont, training on
billions of US quotes, found a stable, universal relation between order-flow history and the direction of the next
move, and that a model pooled across stocks beats stock-specific ones: scale is what made the raw input work. Kolm,
Turiel and Westray, on 115 Nasdaq stocks, found that simpler networks trained on order flow significantly outperform
most models trained directly on order books. The LOBCAST benchmark of fifteen published deep models (Prata et al., 2024)
found that all of them lose a large part of their performance on data they were not tuned on. Two lessons carry over to
any order-book model: build the stationary features the microstructure literature already knows (order flow, imbalance,
Book~7, chapter~8), and judge a model on sessions, days and instruments it has never seen.

\begin{method}[Training a sequence model on order books]\label{met:ml:lobseq:method}
\begin{enumerate}
\item Split by session (day, instrument), never by window; fit every normalisation on training sessions.
\item Fix the label's horizon and dead band from the trade the model serves; report the class shares and a score that
does not reward predicting the majority class (the IC, or a cost-aware P\&L).
\item Start from a linear model on order-flow and imbalance features; give the sequence model the same features as extra
channels before asking it to find them.
\item Report the spread across sessions and seeds, and a learning curve in sessions.
\end{enumerate}
\end{method}

\section{Tutorial: DeepLOB on a laptop}\label{tut:ml:sequence-models-on-order-books:tut}

\begin{tutorial}
\textbf{Goal.} Turn simulated sessions into order-book windows, train four sequence models and a logistic baseline, and
measure the learning curve in sessions. \textbf{End state:} \cref{fig:ml:lobseq:models,fig:ml:lobseq:scaling}.
\begin{tutsteps}
\item \textbf{Windows and labels}, inside each session.
\omcode{code/firm/lobseq/firm_lobseq.py}{57}{73}{Windows of snapshots, smoothed direction labels, and a scaler fitted on training windows.}\label{lst:ml:lobseq:windows}
\item \textbf{The TCN and the attention model.}
\omcode{code/firm/lobseq/firm_lobseq.py}{97}{136}{Causal dilated convolutions and a one-layer transformer.}\label{lst:ml:lobseq:models}
\item \textbf{Run} \texttt{ml\_lobseq.table()}, \texttt{scaling(4)}, \texttt{scaling(8)}, \texttt{scaling(16)} and
\texttt{fig\_lobseq.py}.
\end{tutsteps}
\textbf{What to change next.} Feed the TCN the hand-made features as ten extra channels and see how much of the gap to the
\omterm{def:ml:linear-and-regularised-baselines:logistic}{logistic regression} closes; replace the smoothed label by the sign of the next mid change.
\end{tutorial}

\section{Build: order-book sequence datasets and models}\label{bld:ml:sequence-models-on-order-books:lobseq}

\begin{build}
\bfield{Purpose} Order-book sequence models trained and scored the same way on simulated sessions today and on Book~10's
exchange simulator's recorded tapes when they are available.
\bfield{Interface} \texttt{snapshots(tape, step, levels, start)}, \texttt{windows(S, mid, T, horizon, threshold)},
\texttt{Scaler}, \texttt{CNN}, \texttt{LSTMNet}, \texttt{TCN}, \texttt{Attention}, \texttt{train(model, X, y, Xv, yv,
epochs, lr, batch, patience, seed)}, \texttt{evaluate(model, X, y, fwd)}.
\bfield{Rules} A session's windows stay in one split; the book is rebuilt from messages with nothing after the snapshot's
time; one CPU thread, deterministic algorithms.
\bfield{Acceptance tests} \texttt{code/firm/lobseq/tests/}: snapshots match a book rebuilt by hand at a given time and use
no later message; windows and labels by hand on a short series; the TCN's output at a step does not depend on later
inputs (causality); each model trains on a planted pattern and beats chance.
\bfield{Stretch} Market-by-order inputs (each order's age and queue position); a model pooled across several simulated
instruments, in the spirit of Sirignano and Cont.
\end{build}

\begin{omsources}
\item Z.~Zhang, S.~Zohren and S.~Roberts, ``DeepLOB: deep convolutional neural networks for limit order books'',
\emph{IEEE Transactions on Signal Processing} 67(11), 2019.
\item J.~Sirignano and R.~Cont, ``Universal features of price formation in financial markets: perspectives from deep
learning'', \emph{Quantitative Finance} 19(9), 2019.
\item P.~N.~Kolm, J.~Turiel and N.~Westray, ``Deep order flow imbalance: extracting alpha at multiple horizons from the
limit order book'', \emph{Mathematical Finance} 33(4), 2023.
\item M.~Prata et al., ``LOB-based deep learning models for stock price trend prediction: a benchmark study'',
\emph{Artificial Intelligence Review} 57, 2024.
\item A.~Ntakaris, M.~Magris, J.~Kanniainen, M.~Gabbouj and A.~Iosifidis, ``Benchmark dataset for mid-price forecasting of
limit order book data with machine learning methods'', \emph{Journal of Forecasting} 37(8), 2018.
\item S.~Hochreiter and J.~Schmidhuber, ``Long short-term memory'', \emph{Neural Computation} 9(8), 1997; S.~Bai, J.~Z.~Kolter
and V.~Koltun, ``An empirical evaluation of generic convolutional and recurrent networks for sequence modeling'', arXiv,
2018; A.~Vaswani et al., ``Attention is all you need'', \emph{NeurIPS}, 2017.
\end{omsources}

\section{Exercises}

\begin{exercise}[$\star$]\label{exo:ml:sequence-models-on-order-books:1}
A TCN has filters of width three and dilations 1, 2, 4 and 8. What is its receptive field in snapshots? With width five?
\end{exercise}

\begin{exercise}[$\star$]\label{exo:ml:sequence-models-on-order-books:2}
The test windows are 72.7\% flat, 14.7\% down and 12.5\% up. What accuracy does always predicting ``flat'' reach, and what
does a random guess in the class proportions reach?
\end{exercise}

\begin{exercise}[$\star$]\label{exo:ml:sequence-models-on-order-books:3}
One attention head over $T = 20$ positions of width $h = 32$: how many multiplications does $\mathsf Q\mathsf K^\top$
take, and how does it scale with $T$?
\end{exercise}

\begin{exercise}[$\star\star$]\label{exo:ml:sequence-models-on-order-books:4}
Why does the CNN, which averages over the window, do so badly, and what small change would fix it?
\end{exercise}

\begin{exercise}[$\star\star$]\label{exo:ml:sequence-models-on-order-books:5}
The label averages the mid over the next ten seconds. Why does that make it easier to predict than the next mid change,
and what does it mean for trading?
\end{exercise}

\begin{exercise}[$\star\star$]\label{exo:ml:sequence-models-on-order-books:6}
\emph{Find the flaw.} ``We cut our 50 days of order-book windows into train and \omterm{def:ml:why-financial-machine-learning-is-different:sets}{test sets} at random and our network is 88\%
accurate.''
\end{exercise}

\begin{exercise}[$\star\star\star$]\label{exo:ml:sequence-models-on-order-books:7}
\emph{Coding.} Train the LSTM on sixteen sessions (\texttt{scaling(16, name='LSTM')}). Report its IC and compare with the
TCN's at sixteen sessions.
\end{exercise}

\begin{exercise}[$\star\star\star$]\label{exo:ml:sequence-models-on-order-books:8}
Show that a stack of \omterm{def:ml:sequence-models-on-order-books:tcn}{causal convolutions} of width $k$ with dilations $1, 2,\dots,2^{L-1}$ has a receptive field of
$1 + (k - 1)(2^L - 1)$ positions and that its output at $t$ depends on no input after $t$.
\end{exercise}

\section{Problem: DeepLOB on a Laptop}

\begin{problem}[{Weekend problem --- the raw book against hand-made features}]\label{pb:ml:sequence-models-on-order-books:1}
The chapter's sessions: eight to train, three to stop, eight to test, eight more to grow the \omterm{def:ml:why-financial-machine-learning-is-different:sets}{training set}.

\medskip\noindent\textbf{Part I --- The data.}
\begin{enumerate}
\item What is in one snapshot and one window?
\item How is the label built, and what are the class shares?
\item Why is accuracy a poor score here, and what does the chapter use instead?
\item Why must windows never cross the split between sessions?
\end{enumerate}

\medskip\noindent\textbf{Part II --- The models.}
\begin{enumerate}[resume]
\item What are the four sequence models' receptive fields over the window?
\item What does each score (accuracy and IC)?
\item What does the \omterm{def:ml:linear-and-regularised-baselines:logistic}{logistic regression} on hand-made features score?
\item How large is the spread across sessions and seeds?
\end{enumerate}

\medskip\noindent\textbf{Part III --- Data.}
\begin{enumerate}[resume]
\item After how many \omterm{def:ml:neural-networks-for-noisy-tabular-data:backprop}{epochs} do the models stop, and why?
\item What does the TCN score with four, eight and sixteen training sessions?
\item Does the \omterm{def:ml:linear-and-regularised-baselines:logistic}{logistic regression} gain from more sessions? Why not?
\item What would the curve need to do for the TCN to overtake the baseline?
\end{enumerate}

\medskip\noindent\textbf{Part IV --- The verdict.}
\begin{enumerate}[resume]
\item What did Kolm, Turiel and Westray find about raw books against order flow?
\item What did the LOBCAST benchmark find about published models on new data?
\item What made Sirignano and Cont's raw-input model work?
\item State the \emph{named result}: the best sequence model's IC and accuracy against the \omterm{def:ml:linear-and-regularised-baselines:logistic}{logistic regression} on order-flow
imbalance, with the spread across seeds and sessions.
\item What would you put in production for a ten-second forecast, and what would you try next?
\item How would you turn the IC into a trading decision, given chapter~2's labels?
\item What changes on Book~10's simulator with several instruments?
\item In one sentence: what does a sequence model need to beat a good feature?
\end{enumerate}
\end{problem}

\section{Interview questions}

\begin{interviewq}[$\star$ \iqroles{mle, researcher}]\label{iq:ml:sequence-models-on-order-books:1}
What is a \omterm{def:ml:sequence-models-on-order-books:tcn}{causal convolution}, and why does it matter for a trading model?
\end{interviewq}

\begin{interviewq}[$\star\star$ \iqroles{mle}]\label{iq:ml:sequence-models-on-order-books:2}
Explain self-attention in a few equations. What does it cost in the length of the sequence?
\end{interviewq}

\begin{interviewq}[$\star\star$ \iqroles{researcher}]\label{iq:ml:sequence-models-on-order-books:3}
Your order-book model is 80\% accurate on a three-class mid-price label. What questions do you ask before you are
impressed?
\end{interviewq}

\begin{interviewq}[$\star\star$ \iqroles{researcher, trader}]\label{iq:ml:sequence-models-on-order-books:4}
Raw order-book levels or order-flow features as inputs? Argue from the evidence.
\end{interviewq}

\begin{interviewq}[$\star\star$ \iqroles{mle}]\label{iq:ml:sequence-models-on-order-books:5}
How do you split order-book data for training and testing, and why?
\end{interviewq}

\begin{interviewq}[$\star\star\star$ \iqroles{researcher, mle}]\label{iq:ml:sequence-models-on-order-books:6}
Why might an LSTM underperform a TCN on short windows of order-book data, and when would you expect the opposite?
\end{interviewq}
